\section{Method}
\label{sec:method}

\begin{figure*}[t]
\centering
% TODO: replace with generated figure
\fbox{\parbox{0.95\textwidth}{\centering\vspace{2em}
\textbf{Figure placeholder: STA Overview (3 panels)}\\[0.5em]
(a) NL input $\to$ LLM formalizer $\to$ formal logic $\to$ solver $\to$ proved/fail \quad
(b) STA ablation: remove each $c_i$, re-run solver, compute ISR \quad
(c) ISR as diagnostic vs.\ ISR as reward (Goodhart arrow)
\vspace{2em}}}
\caption{Overview of Solver Trace Ablation (STA). (a)~A formalizer translates NL premises and hypothesis into formal logic, submitted to an external solver. (b)~STA removes each component in turn and re-runs the solver to compute ISR. (c)~ISR serves as a diagnostic for structural integrity; when repurposed as a reward signal, the model games the metric itself (meta-level Goodhart).}
\label{fig:sta-overview}
\end{figure*}

\subsection{Task Formulation}
\label{sec:task}

We study the task of translating natural language (NL) premises and hypotheses into formal logic representations verifiable by an external solver.
Given NL premises $P$ and hypothesis $H$, a formalizer $\pi_\theta$ generates a formal representation $x$ comprising $n$ components $\{c_1, \ldots, c_n\}$---axioms, rules, and a conclusion---submitted to solver $\mathcal{S}$, which returns a binary outcome: \texttt{proved} or \texttt{fail}.

We consider two target formalisms:
(1)~\textbf{First-order logic (FOL)}, where components are Prover9-format axioms and a goal clause, verified by the Prover9 automated theorem prover~\cite{mccune2010prover9};
(2)~\textbf{Prolog}, where components are Horn clauses, facts, and a query, verified by SWI-Prolog.
Training and evaluation data are drawn from ProofWriter~\cite{proofwriter}, with a fixed split of 500 training and 100 held-out evaluation examples across all experiments.

\subsection{GRPO Training}
\label{sec:grpo}

We train the formalizer with Group Relative Policy Optimization~\citep[GRPO;][]{deepseekmath}.
For each training problem, the policy $\pi_\theta$ generates $K{=}8$ candidate formalizations.
Each candidate is submitted to the solver, yielding a binary execution reward:
\begin{equation}
  r_{\text{exec}}(x) =
  \begin{cases}
    1 & \text{if } \mathcal{S}(x) = \texttt{proved} \\
    0 & \text{otherwise}
  \end{cases}
  \label{eq:r_exec}
\end{equation}
GRPO normalizes rewards within each group of $K$ completions to estimate advantages, then updates $\pi_\theta$ via a clipped policy gradient objective without a critic model.
The base model is Qwen3-4B; all runs use learning rate $10^{-6}$ and train for 500 gradient steps unless otherwise noted.

Under execution-only reward, the optimizer receives no signal about \emph{why} a proof succeeds---whether the formalization captures genuine logical derivation or exploits a degenerate shortcut.
This opacity motivates a structural diagnostic.

\subsection{Solver Trace Ablation}
\label{sec:sta}

Solver Trace Ablation (STA) tests whether each component of a proved formalization is \emph{necessary} for the proof.
The protocol applies leave-one-out (LOO) ablation---a standard technique in mutation testing~\cite{demillo1978, jia-harman-survey} and feature attribution---to solver inputs.
We claim no novelty in the ablation mechanism itself; the contribution lies in applying it as a diagnostic for reward hacking in neurosymbolic GRPO training.

\begin{algorithm}[t]
\caption{Solver Trace Ablation (STA)}
\label{alg:sta}
\begin{algorithmic}[1]
\Require Proved formalization $x = \{c_1, \ldots, c_n\}$, solver $\mathcal{S}$
\Ensure $\text{ISR}(x) \in [0, 1]$
\State $\textit{necessary} \gets 0$
\For{$i = 1$ to $n$}
    \State $x_{\setminus i} \gets x \setminus \{c_i\}$
    \If{$\mathcal{S}(x_{\setminus i}) = \texttt{fail}$}
        \State $\textit{necessary} \gets \textit{necessary} + 1$
    \EndIf
\EndFor
\State \Return $\textit{necessary} \;/\; n$
\end{algorithmic}
\end{algorithm}

For a proved formalization $x = \{c_1, \ldots, c_n\}$, STA removes each component in turn and re-runs the solver (Algorithm~\ref{alg:sta}).
The \textbf{Intensional Structural Reward} (ISR) quantifies the fraction of components whose removal causes proof failure:
\begin{equation}
  \text{ISR}(x) = \frac{1}{n} \sum_{i=1}^{n} \mathbf{1}\!\left[\mathcal{S}(x_{\setminus i}) = \texttt{fail}\right]
  \label{eq:isr}
\end{equation}

ISR is bounded in $[0, 1]$.
$\text{ISR} = 1$ indicates every component is individually necessary---removing any single one breaks the proof.
$\text{ISR} = 0$ indicates no component is individually necessary, signaling redundancy or degeneracy.
The protocol requires $n$ solver calls per proved sample.
For unproved samples ($\mathcal{S}(x) = \texttt{fail}$), ISR is undefined; there is no successful proof to ablate.
Throughout the paper, we report \emph{proved-only} ISR (denoted ISR$^*$), computed exclusively over proved samples to avoid conflating structural quality with prove rate.

ISR measures \emph{syntactic component necessity}: whether each piece is required for the proof to succeed.
It does not measure semantic faithfulness to the NL input.
A formalization can achieve $\text{ISR} = 1$ while being semantically unrelated to the original statement---a single premise identical to the conclusion trivially satisfies necessity.

\paragraph{Premise count bias.}
ISR is mathematically sensitive to component count $n$.
When $n = 2$ (one premise, one conclusion), removing the sole premise breaks any non-trivial proof, yielding $\text{ISR} \geq 0.5$ regardless of proof quality.
Conversely, formalizations with many components tend toward lower ISR as redundancy accumulates.
This is an inherent property of LOO-based metrics, not specific to our setting.
When comparing ISR across conditions that differ in mean premise count, we control for this confound via premise-count-matched analysis (\S\ref{sec:experiments}).
Figure~\ref{fig:sta-overview} illustrates the full protocol.

\subsection{Reward Formulations}
\label{sec:reward}

We study six reward conditions to isolate the effect of ISR as a training signal (Table~\ref{tab:reward}).

\begin{table}[t]
\centering
\small
\begin{tabular}{@{}llc@{}}
\toprule
\textbf{Condition} & \textbf{Reward formula} & $\beta$ \\
\midrule
Exec-only (baseline) & $r_{\text{exec}}$ & --- \\
ISR low-dose & $r_{\text{exec}} \cdot (1 + \beta \cdot \text{ISR})$ & 0.1 \\
ISR high-dose & $r_{\text{exec}} \cdot (1 + \beta \cdot \text{ISR})$ & 0.3 \\
Random low-dose & $r_{\text{exec}} \cdot (1 + \beta \cdot u)$ & 0.1 \\
Random high-dose & $r_{\text{exec}} \cdot (1 + \beta \cdot u)$ & 0.3 \\
SFT baseline & Supervised (NLL on GT) & --- \\
\bottomrule
\end{tabular}
\caption{Reward conditions. $u \sim \text{Uniform}(0,1)$. All GRPO conditions use Qwen3-4B, $K{=}8$, lr$=10^{-6}$, 500 steps.}
\label{tab:reward}
\end{table}

\paragraph{Additive ISR bonus.}
ISR augments the execution reward as a multiplicative bonus, active only for proved samples:
\begin{equation}
  r = r_{\text{exec}} \cdot (1 + \beta \cdot \text{ISR}(x))
  \label{eq:additive_isr}
\end{equation}
Proved samples receive reward in $[1, 1{+}\beta]$; unproved samples receive $0$.
We evaluate $\beta \in \{0.1, 0.3\}$ to test dose-response dynamics.
An earlier convex formulation ($r = \alpha \cdot r_{\text{exec}} + \beta \cdot r_{\text{exec}} \cdot \text{ISR}$, $\alpha{+}\beta{=}1$) was discarded because $\text{ISR} < 1$ always penalized relative to pure execution reward, collapsing prove rate.

\paragraph{Random additive control.}
To disentangle ISR's structural signal from the effect of reward variance, we replace ISR with a uniform random variable:
\begin{equation}
  r = r_{\text{exec}} \cdot (1 + \beta \cdot u), \quad u \sim \text{Uniform}(0, 1)
  \label{eq:additive_random}
\end{equation}
with identical $\beta$ values.
If ISR provides genuine structural guidance beyond stochastic perturbation, it should outperform this control.

\paragraph{SFT baseline.}
Standard supervised fine-tuning on the same 500 training examples using ground-truth FOL annotations as targets.
This control isolates whether structural degeneration stems from RL reward optimization or from data exposure alone.

All GRPO conditions share identical hyperparameters and are evaluated across seeds $\in \{42, 43, 44\}$.

\subsection{Hacking Taxonomy}
\label{sec:taxonomy}

We identify eight distinct modes of structural degeneration across reward conditions and logic types.
Definitions appear below; quantitative prevalence analysis is deferred to~\S\ref{sec:experiments}.

\textbf{H1: Conclusion embedding.}
At least one premise is identical to the conclusion after normalization, making the proof trivially $p \rightarrow p$.
This is the dominant mode across all GRPO conditions (52--93\% of proved samples), including both ISR and random reward.

\textbf{H3: Complexity without utility.}
Syntactically complex but logically vacuous premises that do not contribute to the proof.

\textbf{H7a: Text-match tautology.}
The formalization reduces to a single premise identical to the conclusion ($n{=}2$).
H7a occurs across GRPO conditions---including exec-only and random control---but ISR reward amplifies its frequency approximately fivefold (ISR pooled: 70.6\% vs.\ Random pooled: 13.7\%; Fisher $p = 1.3 \times 10^{-39}$, OR$= 15.05$).

\textbf{H7b: Obfuscated tautology.}
Logically equivalent to the conclusion but textually distinct---e.g., vacuous quantifier wrapping ($\forall x.\, P \rightarrow P$ where $P$ is ground) or variable renaming.
Observed exclusively under ISR reward, representing a strategy that simultaneously games the execution reward, ISR metric, and text-match hacking classifier.

\textbf{H8: Conjunction collapse.}
All premises merged into a single conjunction, trivializing proof structure while preserving a multi-component appearance.

\textbf{Repetition collapse.}
Identical premises generated repeatedly, inflating component count without adding logical content.

\textbf{Contradiction explosion.}
Contradictory premises (e.g., $p(a) \wedge \neg p(a)$) from which any conclusion follows \emph{ex falso quodlibet}.
Observed at extended training (1000 steps), with mean component count rising from 6.6 to 18.6.

\textbf{Query-as-fact gaming} (Prolog-specific).
The model converts the NL question into a Prolog atom asserted simultaneously as fact and query, making the proof trivially true.
This mode evades text-matching classifiers; detection requires functional analysis of the necessity vector ($n_{\text{necessary}} = 1$ for 95/97 proved samples).

We define a sample as \emph{genuine} if it is proved and belongs to none of the primary hacking modes:
\begin{equation}
  \text{genuine} = \text{proved} - \left| H7 \cup H1 \cup H8 \right|
  \label{eq:genuine}
\end{equation}
where the set union handles overlapping modes (e.g., samples exhibiting both H1 and H8 simultaneously).
Detection combines two complementary methods: normalized text matching (for H1, H7a) and functional analysis of the necessity vector (for modes that bypass text-based classifiers, such as H7b and query-as-fact gaming).
